\documentclass[12pt]{article}

%% ── Encoding & font ──────────────────────────────────────────────────────────
\usepackage[T1]{fontenc}
\usepackage[utf8]{inputenc}
\usepackage{mathptmx}            % Times New Roman for text and math

%% ── Page layout ─────────────────────────────────────────────────────────────
\usepackage[margin=1in]{geometry}
\usepackage{setspace}
\singlespacing
\setlength{\parskip}{0pt}
\setlength{\parindent}{1.5em}

%% ── Bibliography (APA author-year) ──────────────────────────────────────────
\usepackage{natbib}
\bibliographystyle{apalike}

%% ── Tables ───────────────────────────────────────────────────────────────────
\usepackage{booktabs}
\usepackage{multirow}
\usepackage{array}
\usepackage{float}
\usepackage{colortbl}
\usepackage[labelfont=bf,font=small]{caption}

%% ── Graphics / drawing ───────────────────────────────────────────────────────
\usepackage{graphicx}
\usepackage{tikz}
\usetikzlibrary{shapes.geometric,arrows.meta,positioning,fit,backgrounds,calc}

%% ── Math ─────────────────────────────────────────────────────────────────────
\usepackage{amsmath}
\usepackage{amssymb}

%% ── Colours ──────────────────────────────────────────────────────────────────
\usepackage{xcolor}
\definecolor{dknavy}{HTML}{1E3A5F}
\definecolor{dkblue}{HTML}{2563EB}
\definecolor{dkgreen}{HTML}{16A34A}
\definecolor{dkamber}{HTML}{D97706}
\definecolor{dkcyan}{HTML}{0891B2}
\definecolor{dkpurple}{HTML}{7C3AED}
\definecolor{rowgray}{HTML}{F3F4F6}

%% ── Miscellaneous ────────────────────────────────────────────────────────────
\usepackage{enumitem}
\usepackage[hidelinks]{hyperref}
\usepackage{microtype}

%% ── Section formatting (IEEE-style Roman + letter numbering) ─────────────────
\usepackage{titlesec}
\titleformat{\section}
  {\normalfont\bfseries\centering}
  {\Roman{section}.}{0.5em}{\MakeUppercase}
\titleformat{\subsection}
  {\normalfont\itshape\bfseries}
  {\Alph{subsection}.}{0.5em}{}
\titleformat{\subsubsection}
  {\normalfont\itshape}
  {\arabic{subsubsection})}{0.5em}{}
\titlespacing*{\section}     {0pt}{10pt}{4pt}
\titlespacing*{\subsection}  {0pt}{7pt}{2pt}
\titlespacing*{\subsubsection}{0pt}{5pt}{1pt}

%% ── Running header ───────────────────────────────────────────────────────────
\usepackage{fancyhdr}
\pagestyle{fancy}\fancyhf{}
\lhead{\small\textit{IS596 Final Project}}
\rhead{\small\textit{Deep-Guard Agent}}
\cfoot{\thepage}
\renewcommand{\headrulewidth}{0.4pt}

%% ─────────────────────────────────────────────────────────────────────────────
\begin{document}

%% ── Title block ──────────────────────────────────────────────────────────────
\begin{center}
  {\Large\bfseries Deep-Guard Agent: Bimodal Audio-Visual Deepfake\\[2pt]
  Detection with Forensic Legal Reporting}\\[10pt]
  {\normalsize
    Qiming Li\,$^{1}$\quad Yiting Wang\,$^{1}$\quad Yawen Ou\,$^{1}$}\\[3pt]
  {\small $^{1}$IS596 --- Interactive Intelligent Systems, Final Project}
\end{center}

\medskip

%% ── Abstract ─────────────────────────────────────────────────────────────────
\noindent\textbf{\textit{Abstract}}\,---\,
The proliferation of generative-AI media has created urgent demand for deepfake
detection tools that are accurate, explainable, and legally actionable.
Existing systems are predominantly unimodal and return opaque probability scores
without supporting evidence.
We present \textbf{Deep-Guard Agent}, a six-module pipeline that exploits
articulatory physics to detect deepfakes: the physical constraints governing
vocal-tract motion during speech are systematically violated by current
generative models, producing measurable audio-visual discrepancies.
The system combines a MediaPipe visual lip encoder with a Wav2Vec 2.0
audio-articulatory encoder, aligns both streams via a Temporal Transformer
fusion module, and synthesises findings through a LLaMA 3.3 70B reasoning layer.
A nine-section ISO/IEC 27037-compliant forensic report with SHA-256
chain-of-custody is generated automatically.
A formative user study ($N\!=\!10$) demonstrates a 30-percentage-point accuracy
improvement over an unaided baseline (85\% vs.\ 55\%) and a mean SUS score of
70.0/100, confirming above-average usability.

\smallskip
\noindent\textbf{\textit{Index Terms}}\,---\,deepfake detection,
audio-visual alignment, articulatory representations, large language models,
forensic reporting, explainable AI.

\medskip

%% ── I. Introduction ──────────────────────────────────────────────────────────
\section{Introduction}

The emergence of high-fidelity generative models has dramatically lowered the
barrier to creating photorealistic synthetic video.
Since 2019, reported deepfake incidents have grown by approximately 3{,}000\%,
with an estimated 96\% targeting individuals non-consensually
\citep{sensity2023}.
The consequences extend well beyond personal harm: deepfakes are weaponised to
fabricate political speech, manufacture false testimony, and erode institutional
trust at scale.
Concurrently, new legislation --- including EU AI Act Article~50
\citep{euaiact2024} and the US Take It Down Act \citep{takedown2025} ---
requires platforms to supply verifiable forensic evidence of synthetic media,
creating urgent demand for documentation-grade detection tools.

Three critical gaps characterise the current detection landscape.
\textit{First}, most production tools are unimodal, examining video \textit{or}
audio in isolation.
This ignores the fundamental physical constraint governing authentic speech:
vocal-tract articulation creates a tightly coupled, physically consistent
co-dependence between lip motion and acoustic output that generative adversarial
networks (GANs) and diffusion-based synthesis pipelines cannot reliably replicate
\citep{wang2024artavdf}.
\textit{Second}, existing tools function as black boxes, returning a score
without any supporting rationale.
This opacity prevents users from developing calibrated trust and forecloses the
cognitive scaffolding that enables critical media literacy.
\textit{Third}, no open system produces forensic-grade documentation satisfying
US federal evidentiary requirements under the Daubert standard
\citep{daubert1993} or international digital-evidence guidelines under
ISO/IEC~27037:2012 \citep{iso27037}.

This paper introduces \textbf{Deep-Guard Agent}, which closes all three gaps
simultaneously through three primary contributions:
\begin{enumerate}[leftmargin=0.9in, noitemsep, topsep=2pt]
  \item \textbf{Bimodal Detection Architecture.}
    A Temporal Transformer fusion module that aligns Wav2Vec 2.0 audio embeddings
    \citep{baevski2020} with MediaPipe lip feature vectors and produces per-frame
    discrepancy scores grounded in articulatory physics.
  \item \textbf{LLM-Powered Explainability.}
    A JSON-structured reasoning pipeline using LLaMA 3.3 70B that converts
    numerical detection evidence into chain-of-thought verdicts with supporting
    rationale and recommended actions.
  \item \textbf{Legal Forensic Report Generator.}
    An HTML report module producing SHA-256-hashed, chain-of-custody
    documentation conforming to ISO/IEC~27037:2012 \citep{iso27037} and the
    Daubert admissibility standard \citep{daubert1993}.
\end{enumerate}

%% ── II. Related Work ─────────────────────────────────────────────────────────
\section{Related Work}

\subsection{Audio-Visual Speech Representations}

The theoretical basis for bimodal deepfake detection derives from self-supervised
audio-visual speech learning.
\citet{shi2022avhubert} introduced AV-HuBERT, demonstrating that joint
audio-visual representation learning outperforms unimodal approaches on
lip-reading and speech recognition benchmarks.
Their work established that temporal audio-visual alignment is a learnable signal
with generalisable discriminative power.
Complementarily, \citet{baevski2020} proposed Wav2Vec 2.0, a contrastive
pre-training framework mapping raw waveforms to high-level articulatory
representations; trained on 960 hours of LibriSpeech, its latent space captures
the physical mechanics of vocal production, making it sensitive to the synthetic
artefacts introduced when generative pipelines decouple audio from authentic lip
dynamics.

\subsection{Deepfake Detection Methods}

\citet{wang2024artavdf} provided the closest methodological precedent,
demonstrating that articulatory representation learning significantly improves
detection across the FaceForensics++ benchmark \citep{rossler2019}.
Their ART-AVDF framework established that phoneme-level inconsistencies ---
particularly violations of bilabial closure during voiced plosives /b/, /p/, and
/m/ --- are reliable forgery indicators.
Prior spatial-artefact approaches \citep{rossler2019} have proven increasingly
brittle as synthesis quality improves and pixel-level GAN artefacts diminish;
cross-modal inconsistency provides a more principled detection surface because it
is grounded in physical law rather than distributional statistics.

\subsection{Explainability and Human Decision Support}

The cognitive dimension of deepfake harm has received insufficient attention in
the detection literature.
Providing a bare probability score can paradoxically increase misplaced
confidence by anchoring users to an authority signal they cannot interrogate
\citep{abercrombie2024}.
Our LLM reasoning layer is designed to shift users from intuitive System~1
processing toward analytical System~2 evaluation \citep{kahneman2011}, a
mechanism that has been shown to reduce susceptibility to misinformation when
effectively scaffolded.

\subsection{Forensic and Legal Standards}

The Daubert standard requires scientific evidence to be testable, peer-reviewed,
associated with a known error rate, and generally accepted in the relevant
community \citep{daubert1993}.
ISO/IEC~27037:2012 extends these requirements to digital forensic evidence,
specifying chain-of-custody documentation and integrity verification
\citep{iso27037}.
To our knowledge, Deep-Guard Agent is the first open detection tool to
operationalise both standards through a programmatically generated report.

%% ── III. System Architecture ─────────────────────────────────────────────────
\section{System Architecture}

\subsection{Overview}

Deep-Guard Agent consists of six modules in a two-stage pipeline
(Figure~\ref{fig:arch}).
The \textit{encoding stage} (Modules 1--3) independently extracts multimodal
representations from raw video input.
The \textit{reasoning stage} (Modules 4--6) fuses these representations,
generates a natural-language analysis, and produces structured output artefacts.

\subsection{Visual Lip Encoder}

Each decoded video frame is processed by MediaPipe FaceLandmarker to localise
478 facial landmarks at sub-pixel resolution.
Six lip-specific descriptors are computed per frame: vertical openness,
aspect ratio, horizontal asymmetry, elliptical eccentricity, corner angle, and a
binary bilabial closure flag.
Concatenated with raw coordinate values for 20 outer and inner lip keypoints,
these form a 256-dimensional feature vector
$\mathbf{v}_{t} \in \mathbb{R}^{256}$ at each timestep $t$.
A weighted moving average ($\alpha = 0.7$) suppresses high-frequency tracking
jitter while preserving articulatory dynamics.

\subsection{Audio-Articulatory Encoder}

The audio stream is extracted via FFmpeg at 16 kHz mono and passed through the
\texttt{facebook/wav2vec2-base-960h} transformer encoder.
The model produces a 768-dimensional hidden-state representation
$\mathbf{a}_{t} \in \mathbb{R}^{768}$ at approximately 50 frames per second,
temporally aligned to the video frame rate through linear interpolation.
Because Wav2Vec 2.0 was pre-trained on authentic human speech, its latent space
encodes articulatory co-production constraints --- particularly the coarticulation
dynamics of voiced bilabial plosives --- that current generative models
systematically violate.

\subsection{Cross-Modal Fusion}

The fusion module projects audio and visual features into a shared
$d = 256$-dimensional space via learned linear projections
$W_{a} \in \mathbb{R}^{768 \times 256}$ and
$W_{v} \in \mathbb{R}^{256 \times 256}$.
A two-layer Temporal Attention encoder with 4-head multi-head self-attention
captures multi-scale phoneme transition dynamics across a 10-frame context
window.
The per-frame discrepancy score is computed as:

\begin{equation}
  s_{t} = \alpha\, s_{t}^{\,\text{learned}}
        + (1-\alpha)\, s_{t}^{\,\text{cosine}}, \quad \alpha = 0.7
  \label{eq:score}
\end{equation}

where $s_{t}^{\,\text{learned}}$ is the sigmoid output of a binary
classification head applied to the concatenated projected features, and
$s_{t}^{\,\text{cosine}} = 1 - \cos(\hat{\mathbf{a}}_{t}, \hat{\mathbf{v}}_{t})$
is the normalised cosine discrepancy.
Frames with $s_{t} \geq 0.65$ are flagged as anomalous.
The video-level score averages per-frame scores, weighted by the flagged-frame
ratio.

\subsection{LLM Reasoner}

The LLM reasoner receives the fusion result --- aggregate score, flagged-frame
set, mean cosine similarity, and flagged ratio --- and issues a structured prompt
to LLaMA 3.3 70B via the Groq inference API (sub-second time-to-first-token).
The prompt enforces a chain-of-thought schema requiring the model to:
(i) enumerate supporting evidence items,
(ii) assign a categorical verdict
    (\textit{Authentic}, \textit{Likely Fake}, or \textit{Suspicious}),
(iii) estimate confidence, and
(iv) recommend follow-up actions.
The response is parsed into a typed JSON object; a rule-based fallback provides
graceful degradation when the API is unavailable, ensuring outputs are always
grounded in the quantitative fusion evidence rather than unconstrained generation.

\subsection{Output Module}

Four artefacts are produced per analysis run:
(i)~an annotated MP4 with per-frame face bounding boxes, lip contour overlays,
and a verdict badge;
(ii)~an interactive HTML analysis report with an inline SVG discrepancy timeline;
(iii)~a nine-section legal forensic report; and
(iv)~a machine-readable JSON file for downstream integration.
UUID-prefixed filenames prevent session collisions under concurrent access.
All LLM-generated strings are sanitised via Python's \texttt{html.escape()}
before template injection, preventing cross-site scripting vulnerabilities.

%% ── IV. Implementation ───────────────────────────────────────────────────────
\section{Implementation}

\subsection{Technology Stack}

Deep-Guard Agent is implemented in Python~3.13.
Core dependencies include MediaPipe 0.10 (visual encoding),
HuggingFace Transformers 4.x with \texttt{facebook/wav2vec2-base-960h}
(audio encoding), PyTorch 2.x (fusion module), Groq SDK (LLM inference),
OpenCV and FFmpeg (video I/O), and Gradio 4.x (web interface).
The complete source code is publicly released as an open-source repository.

\subsection{Engineering Robustness}

Four design decisions ensure production-grade reliability.

\textit{Thread safety.}
The singleton pipeline is guarded by a double-checked locking pattern using
\texttt{threading.Lock()}, ensuring that only one initialisation occurs under
concurrent Gradio requests and preventing race conditions in lazy-loading.

\textit{Efficient frame annotation.}
Flagged frame indices are materialised into a Python \texttt{set} before video
rendering, reducing per-frame membership testing from O($n$) list scans to
O(1) hash lookups --- a factor-of-$n$ speedup for long-form video.

\textit{XSS prevention.}
All LLM-generated strings are passed through \texttt{html.escape()} before HTML
template injection, neutralising prompt-injection payloads that downstream report
viewers could otherwise execute.

\textit{Pure SVG visualisations.}
All charts are rendered as server-side SVG strings, circumventing Gradio's
JavaScript sanitisation layer that strips \texttt{<script>} elements from
rendered HTML components, ensuring charts display correctly without any
client-side dependencies.

\subsection{Legal Forensic Report Design}

The legal forensic report is structured into nine sections across three
functional groups (Table~\ref{tab:legal}).
Each report embeds a SHA-256 hash of the source video file computed at
analysis time, establishing a cryptographic integrity check that verifies
post-submission file authenticity.
The report explicitly cites the EU AI Act Article~50 \citep{euaiact2024},
Take It Down Act \citep{takedown2025}, NIST SP~800-86 \citep{nist80086},
ISO/IEC~27037:2012 \citep{iso27037}, and the Daubert standard \citep{daubert1993},
situating each finding within its applicable legal framework.

%% ── V. Evaluation ────────────────────────────────────────────────────────────
\section{Evaluation}

\subsection{Study Design}

A within-subjects formative user study was conducted with $N = 10$ university
students (5 male, 5 female; age 18--25).
Participants were recruited across two background groups: technical
($n = 5$; Computer Science or Information Systems) and non-technical
($n = 5$; Law, Journalism, Social Science, Psychology).
One participant had prior experience with a deepfake detection tool;
mean AI familiarity was 4.3/7 (SD = 1.6).

Each 60-minute session followed a standardised protocol.
Participants completed a pre-questionnaire assessing demographics and technology
familiarity, followed by an unaided two-video baseline task measuring unassisted
detection ability.
They then used Deep-Guard Agent to analyse three task videos of increasing
ambiguity: V1 (authentic), V2 (clearly manipulated), and V3 (subtle
manipulation).
Post-task instruments included the 10-item System Usability Scale
\citep{brooke1996} and four custom 7-point Likert scales measuring report
comprehension, legal report usefulness, trust in the AI verdict, and overall
satisfaction.
Sessions concluded with a 10-minute semi-structured interview; responses were
inductively coded into ten thematic categories.

\subsection{Quantitative Results}

\textit{Usability.}
The mean SUS score was 70.0 (SD = 13.6), placing the system in the
``Above Average'' category on the \citet{bangor2009} adjective scale.
Five of ten participants exceeded the industry benchmark of 68 \citep{brooke1996}.
A marked background effect was observed: technical participants averaged 83.8
versus 60.8 for non-technical participants, a gap of 22 SUS points
(Table~\ref{tab:results}).

\textit{Detection accuracy.}
Unaided baseline accuracy was 55\%, near chance level, consistent with findings
that deepfakes are visually indistinguishable for untrained observers
\citep{westerlund2019}.
System-assisted accuracy rose to 85\% --- a 30 percentage point improvement.
Agreement rates varied by stimulus: 100\% for V1 (authentic),
95\% for V2 (clearly fake), and 60\% for V3 (subtle), confirming that ambiguous
cases remain challenging even with AI assistance.

\textit{Subjective scales.}
All four Likert constructs fell in the neutral-to-positive range (4.17--4.72/7).
Trust in the AI verdict was rated highest ($M = 4.72$, SD = 0.34).
Report comprehension was rated lowest ($M = 4.17$, SD = 1.17), identifying
technical vocabulary as the primary usability barrier.
Legal report usefulness was rated significantly higher by law and journalism
participants ($M = 5.75$) than by CS/IS participants ($M = 4.62$), affirming the
report's relevance to its intended professional audience.

\textit{Correlation.}
A Spearman rank correlation between AI familiarity and SUS score yielded
$\rho = 0.97$ ($p < .001$), indicating a near-perfect monotonic association.
Deepfake familiarity showed negligible correlation with trust ($\rho = 0.19$),
suggesting that general AI experience --- rather than domain-specific knowledge
--- is the primary driver of usability perception.

\subsection{Qualitative Themes}

Inductive coding of interview transcripts identified four dominant themes.
Positive themes included \textit{Clear Verdict} (7/10 participants) and
\textit{Trust in AI Verdict} (6/10), indicating that the verdict badge and
confidence score communicated effectively to most users.
Six participants affirmed the practical value of the legal report, with the Law
student noting they could envisage submitting the document as formal evidence ---
directly validating the system's novel contribution.
Concern themes centred on \textit{Confusing Technical Terminology} (4/10)
--- terms such as \textit{cosine similarity} and
\textit{articulatory discrepancy} were opaque to non-technical users ---
\textit{Processing Speed} (4/10), and \textit{Unclear Timeline Annotations}
(4/10).

%% ── VI. Discussion ───────────────────────────────────────────────────────────
\section{Discussion}

\subsection{Reliability of LLM-Generated Explanations}

A pertinent question concerns whether the LLM reasoner introduces hallucinations
or inaccurate interpretations.
The system employs two structural defences.
First, the LLM receives exclusively quantitative, structured inputs (aggregate
score, flagged-frame ratio, mean cosine similarity) rather than raw media,
constraining the generative space to numerically grounded interpretations.
Second, the JSON output schema enforces typed fields with enumerated verdict
values, preventing free-form speculation.
Nevertheless, we acknowledge the absence of an explicit faithfulness
verification layer as a current limitation: the system does not cross-validate
generated evidence items against input statistics using a natural language
inference model.
A natural extension would apply lightweight NLI-based consistency checking
\citep{abercrombie2024} to each generated evidence item before report rendering.

\subsection{Limitations}

Three limitations bound the current scope.
\textit{Untrained fusion module.}
The fusion module operates with randomly initialised weights; no labelled
deepfake dataset was used for supervised training.
Discrepancy scores therefore reflect relative articulatory patterns rather than
calibrated posterior probabilities.
Detection performance is not yet comparable to supervised systems trained on
FaceForensics++ \citep{rossler2019}; however, the architecture is checkpoint-ready
and this remains the highest-priority development target.
\textit{Sample size.}
The formative study ($N = 10$) is appropriate for theme identification but
underpowered for inferential testing.
\textit{Processing latency.}
Analysing a one-minute video requires approximately five minutes on CPU hardware,
constraining use to asynchronous review workflows.

\subsection{Future Work}

Three development priorities follow directly from the evaluation findings.
Supervised training of the fusion module on FaceForensics++ \citep{rossler2019}
or DFDC would calibrate discrepancy scores and enable benchmarked classification
performance.
Addressing the vocabulary gap through inline glossaries and plain-language
summary panels would reduce the 22-point SUS gap between technical and
non-technical users identified in the study.
Streaming inference architecture would extend the system to near-real-time
live-broadcast moderation, the next critical deployment context as regulatory
obligations expand across jurisdictions.

%% ── VII. Conclusion ──────────────────────────────────────────────────────────
\section{Conclusion}

This paper presented Deep-Guard Agent, a six-module open pipeline integrating
bimodal audio-visual deepfake detection, LLM-powered explainability, and
ISO/IEC 27037-compliant forensic documentation in a single system.
Grounded in articulatory physics and self-supervised speech representations
\citep{baevski2020,shi2022avhubert,wang2024artavdf}, the system exploits the
physical constraints governing vocal-tract motion as a detection signal that
current generative models cannot fully replicate.

A formative user study confirmed a 30 percentage point accuracy improvement over
unaided detection and above-average system usability (SUS = 70.0).
The legal forensic report was rated most valuable by law and journalism
participants --- the primary intended users --- validating the system's potential
as a digital-advocacy and evidence-gathering tool.
The principal limitation, an untrained fusion module, is a tractable engineering
challenge rather than a fundamental design constraint.
With supervised training, the architecture is positioned to deliver calibrated,
explainable deepfake detection that is simultaneously accessible to lay users
and admissible as forensic evidence under prevailing legal standards.

%% ── References ───────────────────────────────────────────────────────────────
\bibliography{references}

%% ── Figure: System Architecture ─────────────────────────────────────────────
\clearpage
\begin{figure}[H]
\centering
\begin{tikzpicture}[
  node distance = 0.5cm and 0.65cm,
  mod/.style    = {rectangle, rounded corners=4pt,
                   minimum width=3.0cm, minimum height=1.1cm,
                   text width=2.85cm, align=center,
                   font=\small\bfseries},
  sublabel/.style = {font=\scriptsize, align=center},
  arr/.style      = {-Stealth, thick, color=gray!60},
  darr/.style     = {-Stealth, thick, color=gray!50},
]
%% --- Row 1 (encoding stage) ---
\node[mod, fill=dknavy!85, text=white]  (n1) {M1: Video Input\\[1pt]
  {\scriptsize MP4 / AVI / MOV}};
\node[mod, fill=dkblue!80, text=white,
      right=of n1]                      (n2) {M2: Visual Encoder\\[1pt]
  {\scriptsize MediaPipe \textbullet{} 256-d}};
\node[mod, fill=dkpurple!75, text=white,
      right=of n2]                      (n3) {M3: Audio Encoder\\[1pt]
  {\scriptsize Wav2Vec 2.0 \textbullet{} 768-d}};

%% --- Row 2 (reasoning stage) ---
\node[mod, fill=dkcyan!75, text=white,
      below=1.3cm of n1]                (n4) {M4: Cross-Modal Fusion\\[1pt]
  {\scriptsize Temporal Transformer}};
\node[mod, fill=dkamber!85, text=black,
      below=1.3cm of n2]                (n5) {M5: LLM Reasoner\\[1pt]
  {\scriptsize LLaMA 3.3 70B (Groq)}};
\node[mod, fill=dkgreen!70, text=black,
      below=1.3cm of n3]                (n6) {M6: Output\\[1pt]
  {\scriptsize Report \textbullet{} Video \textbullet{} JSON}};

%% --- Horizontal arrows row 1 ---
\draw[arr] (n1) -- (n2);
\draw[arr] (n2) -- (n3);
%% --- Down arrows ---
\draw[darr] (n1.south) -- (n4.north);
\draw[darr] (n2.south) -- (n5.north);
\draw[darr] (n3.south) -- (n6.north);
%% --- Horizontal arrows row 2 ---
\draw[arr] (n4) -- (n5);
\draw[arr] (n5) -- (n6);

%% --- Stage background boxes ---
\begin{pgfonlayer}{background}
  \node[fill=dkblue!6, rounded corners=5pt,
        inner sep=8pt,
        fit=(n1)(n2)(n3),
        label={[font=\footnotesize\itshape,
                anchor=south west]south west:Encoding Stage}] {};
  \node[fill=dkgreen!6, rounded corners=5pt,
        inner sep=8pt,
        fit=(n4)(n5)(n6),
        label={[font=\footnotesize\itshape,
                anchor=north west]north west:Reasoning Stage}] {};
\end{pgfonlayer}

%% --- Fusion equation note ---
\node[below=0.3cm of n4,
      font=\scriptsize\itshape, text=gray!70]
  {$s_{t}=0.7\,s_{t}^{\text{learned}}+0.3\,s_{t}^{\text{cosine}}$;\;
   flag if $s_{t}\ge 0.65$};
\end{tikzpicture}
\caption{Deep-Guard Agent six-module pipeline. The encoding stage (top)
independently extracts 256-d visual lip features and 768-d audio-articulatory
features. The reasoning stage (bottom) fuses them via Equation~\ref{eq:score},
generates a chain-of-thought LLM verdict, and renders four output artefacts.}
\label{fig:arch}
\end{figure}

%% ── Table 1: Evaluation Results ──────────────────────────────────────────────
\begin{table}[H]
\centering
\caption{Summary of user-study results ($N=10$). SUS grading follows
\citet{bangor2009}; Likert items scored 1--7.}
\label{tab:results}
\small
\renewcommand{\arraystretch}{1.15}
\begin{tabular}{p{5.8cm} c c}
\toprule
\textbf{Measure} & \textbf{Value} & \textbf{Interpretation} \\
\midrule
\multicolumn{3}{l}{\cellcolor{rowgray}\textit{\textbf{System Usability Scale}}} \\
\quad Mean SUS score             & 70.0 (SD\,=\,13.6)  & Above Average (C) \\
\quad Participants $\ge$ 68      & 5/10\,(50\%)         & Above benchmark \\
\quad Technical background       & $M$\,=\,83.8         & Good (B) \\
\quad Non-technical background   & $M$\,=\,60.8         & Marginal (D) \\
\midrule
\multicolumn{3}{l}{\cellcolor{rowgray}\textit{\textbf{Detection Accuracy}}} \\
\quad Unaided baseline           & 55\%                 & Near chance \\
\quad System-assisted (overall)  & 85\%                 & +30\,pp improvement \\
\quad V1 Authentic               & 100\%                & --- \\
\quad V2 Likely Fake             & 95\%                 & --- \\
\quad V3 Suspicious (subtle)     & 60\%                 & Ambiguous cases harder \\
\midrule
\multicolumn{3}{l}{\cellcolor{rowgray}\textit{\textbf{Likert Scales (1--7)}}} \\
\quad Report Comprehension       & 4.17 (SD\,=\,1.17)  & Neutral \\
\quad Legal Report Usefulness    & 4.60 (SD\,=\,0.85)  & Neutral--Positive \\
\quad Trust in AI Verdict        & 4.72 (SD\,=\,0.34)  & Neutral--Positive \\
\quad Overall Satisfaction       & 4.65 (SD\,=\,0.78)  & Neutral--Positive \\
\midrule
\multicolumn{3}{l}{\cellcolor{rowgray}\textit{\textbf{Spearman Correlation}}} \\
\quad AI Familiarity vs.\ SUS    & $\rho=0.97$, $p<.001$& Near-perfect \\
\quad DF Familiarity vs.\ Trust  & $\rho=0.19$          & Negligible \\
\bottomrule
\end{tabular}
\end{table}

%% ── Table 2: Legal Report Structure ─────────────────────────────────────────
\begin{table}[H]
\centering
\caption{Nine-section structure of the legal forensic report and compliance
mapping to applicable standards.}
\label{tab:legal}
\small
\renewcommand{\arraystretch}{1.15}
\begin{tabular}{l l p{5.0cm}}
\toprule
\textbf{Group} & \textbf{Section} & \textbf{Compliance Mapping} \\
\midrule
\multirow{3}{*}{Identification}
  & (1) Case Information  & ISO/IEC 27037 Sec.~5.1 (documentation) \\
  & (2) Forensic Verdict  & Daubert (testability, error rate) \\
  & (3) Methodology       & Daubert (peer-reviewed method) \\
\midrule
\multirow{3}{*}{Evidence}
  & (4) Quantitative Findings & EU AI Act Art.~50 (transparency) \\
  & (5) Discrepancy Timeline  & NIST SP 800-86 (artefact logging) \\
  & (6) Flagged Frames Table  & NIST OpenMFC (temporal evidence) \\
\midrule
\multirow{3}{*}{Compliance}
  & (7) Evidence Summary      & Daubert (general acceptance) \\
  & (8) Chain of Custody      & ISO/IEC 27037 Sec.~6 (integrity) \\
  & (9) Limitations           & ISO/IEC 27037 Sec.~5.4 (limitations) \\
\bottomrule
\end{tabular}
\end{table}

\end{document}
